% Chunk: remainder of \S6.1 Derivation of the Rules (no repeated subsection heading).
% Source: Weinberg QFT Vol. I, physical PDF pp. 289--297 (printed pp. 266--274),
%         beginning with the final paragraph on p. 289 and ending before the \S6.2 heading.
% Status: source-reviewed and compile-clean; equations (6.1.17)--(6.1.31),
%         plus one unnumbered display.
% Figures: 6.3--6.8 integrated at their exact positions with captions.
% Footnotes: 1. Uncertainties: none.

Other combinatoric factors arise when some of the permutations of vertices
have no effect on the Feynman diagram. We noted earlier that the factor $1/N!$
in the series \eqref{eq:6.1.1} is usually cancelled by the sum over the $N!$ diagrams
that differ only in the labelling of the $N$ vertices. However, this
cancellation is incomplete when relabelling the vertices does not yield a new
diagram. This happens most commonly in the calculation of vacuum-to-vacuum
$S$-matrix elements in a theory with a quadratic interaction
$\mathcal H_I=\psi_\ell^\dagger M_{\ell\ell'}\psi_{\ell'}$, where $M$ may
depend on external fields. (The physical significance of such vacuum
fluctuation diagrams is discussed in detail in Volume II.) The Feynman diagram
of $N$th order in $\mathcal H_I$ is a ring with $N$ corners. (See Figure~\ref{fig:6.3}.)
There are only $(N-1)!$ different diagrams here because a permutation of labels
that moves each label to the next vertex around the ring yields the same
diagram. Hence such a graph is accompanied with a factor
\begin{equation}
\frac{(N-1)!}{N!}=\frac{1}{N}.
\tag{6.1.17}
\label{eq:6.1.17}
\end{equation}

\begingroup
\renewcommand{\thefigure}{6.\arabic{figure}}
\setcounter{figure}{2}
\begin{figure}[htbp]
% Figure 6.3 diagram only; caption remains in the chapter text.
\begin{center}
\begin{tikzpicture}[
  x=1cm,
  y=1cm,
  line cap=round,
  line join=round,
  fermion/.style={line width=0.8pt},
  flow/.style={-{Stealth[length=2.5mm,width=1.9mm]},line width=0.8pt},
  external/.style={
    line width=0.8pt,
    decorate,
    decoration={snake,amplitude=1.9pt,segment length=10pt}
  }
]
\coordinate (TL) at (-0.72,1.72);
\coordinate (TR) at (0.72,1.72);
\coordinate (UR) at (1.72,0.72);
\coordinate (LR) at (1.72,-0.72);
\coordinate (BR) at (0.72,-1.72);
\coordinate (BL) at (-0.72,-1.72);
\coordinate (LL) at (-1.72,-0.72);
\coordinate (UL) at (-1.72,0.72);

% Counterclockwise fermion ring.
\draw[fermion] (TR) -- (TL);
\draw[flow] (0.25,1.72) -- (-0.25,1.72);
\draw[fermion] (TL) -- (UL);
\draw[flow] (-1.05,1.39) -- (-1.40,1.04);
\draw[fermion] (UL) -- (LL);
\draw[flow] (-1.72,0.25) -- (-1.72,-0.25);
\draw[fermion] (LL) -- (BL);
\draw[flow] (-1.40,-1.04) -- (-1.05,-1.39);
\draw[fermion] (BL) -- (BR);
\draw[flow] (-0.25,-1.72) -- (0.25,-1.72);
\draw[fermion] (BR) -- (LR);
\draw[flow] (1.05,-1.39) -- (1.40,-1.04);
\draw[fermion] (LR) -- (UR);
\draw[flow] (1.72,-0.25) -- (1.72,0.25);
\draw[fermion] (UR) -- (TR);
\draw[flow] (1.40,1.04) -- (1.05,1.39);

% One external-field line at each interaction vertex.
\draw[external] (TL) -- (-0.95,3.02);
\draw[external] (TR) -- (0.95,3.02);
\draw[external] (UR) -- (3.02,1.02);
\draw[external] (LR) -- (3.02,-1.02);
\draw[external] (BR) -- (0.95,-3.02);
\draw[external] (BL) -- (-0.95,-3.02);
\draw[external] (LL) -- (-3.02,-1.02);
\draw[external] (UL) -- (-3.02,1.02);
\end{tikzpicture}
\end{center}

\caption{An eighth-order graph for the vacuum-to-vacuum amplitude with
particles interacting only with an external field. In this diagram the
external field is represented by wiggly lines. There are $7!$ such diagrams,
differing only by relabelling the vertices, and not counting as different those
labellings that simply rotate the ring. The factor $1/8!$ from the Dyson
formula \eqref{eq:6.1.1} is therefore not entirely cancelled here, leaving us
with an extra factor $1/8$.}
\label{fig:6.3}
\end{figure}
\endgroup

\noindent\textbf{(vi)} In theories involving fermion fields, the use of Eqs.
\eqref{eq:6.1.4}--\eqref{eq:6.1.6} to move annihilation and creation operators to the right and
left introduces minus signs into the contribution of various pairings. To be
specific, we get a minus sign wherever the permutation of the operators in Eq.
\eqref{eq:6.1.1} that is required to put all paired operators adjacent to one another
(with annihilation operators just to the left of the paired creation operators)
involves an odd number of interchanges of fermion operators. (This is because
to compute the contribution of a certain pairing, we can first permute all
operators in Eq. \eqref{eq:6.1.1} so that each annihilation operator is just to the left
of the creation operator with which it is paired, ignoring all commutators and
anticommutators of unpaired operators, and then replace each product of paired
operators with their commutators or anticommutators.) One immediate consequence
is to produce the minus sign in the relative sign of the two terms in Eq.
\eqref{eq:6.1.14} for the fermion propagator. Whatever permutation puts the annihilation
part $\psi^{(+)}(x)$ of a field in $\mathcal H_I(x)$ just to the left of the
creation part $\psi^{(+)\dagger}(y)$ of a field adjoint in $\mathcal H_I(y)$,
the permutation that puts the annihilation part $\psi^{(-)\dagger}(y)$ of the
field adjoint just to the left of the creation part $\psi^{(-)}(x)$ of the field
involves one extra interchange of fermion operators, yielding the minus sign in
the second term of Eq. \eqref{eq:6.1.14} for fermions.

In addition, minus signs can arise in the contribution of whole Feynman
diagrams. As an example, let us take up a theory in which the sole interaction
of fermions takes the form
\begin{equation}
\mathcal H_I(x)=\sum_{\ell m k}g_{\ell m k}\psi_\ell^\dagger(x)
 \psi_m(x)\phi_k(x),
\tag{6.1.18}
\label{eq:6.1.18}
\end{equation}
where $g_{\ell m k}$ are general constants, $\psi_\ell(x)$ are a set of complex
fermion fields, and $\phi_m(x)$ are a set of real bosonic (but not necessarily
scalar) fields. (Not only quantum electrodynamics, but the whole `standard
model' of weak, electromagnetic, and strong interactions, has fermionic
interactions that can all be put in this form.) Let us first take up the
process of fermion--fermion scattering, $12\to1'2'$, to second order in
$\mathcal H_I$. The fermion operators in the second-order term in Eq. \eqref{eq:6.1.1}
appear in the order (with obvious abbreviations)
\begin{equation}
a_{2'}a_{1'}\psi^\dagger(x)\psi(x)\psi^\dagger(y)\psi(y)
 a^\dagger_{1}a^\dagger_{2}.
\tag{6.1.19}
\label{eq:6.1.19}
\end{equation}
There are two connected diagrams to this order, corresponding to the pairings
\begin{equation}
[a_{2'}\psi^\dagger(x)]\,[a_{1'}\psi^\dagger(y)]
 [\psi(y)a^\dagger_{1}]\,[\psi(x)a^\dagger_{2}]
\tag{6.1.20}
\label{eq:6.1.20}
\end{equation}
and
\begin{equation}
[a_{1'}\psi^\dagger(x)]\,[a_{2'}\psi^\dagger(y)]
 [\psi(y)a^\dagger_{1}]\,[\psi(x)a^\dagger_{2}].
\tag{6.1.21}
\label{eq:6.1.21}
\end{equation}
(See Figure~\ref{fig:6.4}.) To go from \eqref{eq:6.1.19} to \eqref{eq:6.1.20} requires an even permutation
of fermionic operators. (For instance, move $\psi(x)$ past three operators to
the right, and then move $a_{1'}$ past one operator to the right.) Thus there is
no extra minus sign in the contribution of the pairing \eqref{eq:6.1.20}. This in itself
is not so important; the \emph{overall} sign of the $S$-matrix does not matter
in transition rates, and in any case depends on sign conventions for the
initial and final states. What is important is that the contributions of
pairings \eqref{eq:6.1.20} and \eqref{eq:6.1.21} have opposite sign, as can be seen most easily
by noting that the only difference between these two pairings is the
interchange of two fermionic operators, $a_{1'}$ and $a_{2'}$. In fact, this
relative minus sign is just what is required by Fermi statistics: it makes the
scattering amplitude antisymmetric under the interchange of particles $1'$ and
$2'$ (or 1 and 2).

\begingroup
\renewcommand{\thefigure}{6.\arabic{figure}}
\setcounter{figure}{3}
\begin{figure}[htbp]
% Figure 6.4 diagram fragment. Caption remains in sec61b.tex.
% Source: physical PDF p. 292 (printed p. 269).
\begin{center}
\begin{tikzpicture}[
  fermion/.style={draw=black, line width=0.8pt},
  boson/.style={draw=black, line width=0.75pt,
    dash pattern=on 3pt off 2.5pt},
  flow/.style={-{Stealth[length=2.4mm,width=1.8mm]}, line width=0.8pt},
  label/.style={font=\small, inner sep=1pt}
]
  % First diagram: 1 and 2 retain their final-state assignments.
  \coordinate (aL) at (-4.35,0);
  \coordinate (aR) at (-1.65,0);
  \coordinate (aUL) at (-5.35,1.55);
  \coordinate (aLL) at (-5.25,-1.45);
  \coordinate (aUR) at (-0.65,1.55);
  \coordinate (aLR) at (-0.75,-1.45);

  \draw[boson] (aL)--(aR);
  \draw[fermion] (aL)--(aUL);
  \draw[fermion] (aLL)--(aL);
  \draw[fermion] (aR)--(aUR);
  \draw[fermion] (aLR)--(aR);
  \draw[flow] ($(aL)!0.37!(aUL)$)--($(aL)!0.62!(aUL)$);
  \draw[flow] ($(aLL)!0.37!(aL)$)--($(aLL)!0.62!(aL)$);
  \draw[flow] ($(aR)!0.37!(aUR)$)--($(aR)!0.62!(aUR)$);
  \draw[flow] ($(aLR)!0.37!(aR)$)--($(aLR)!0.62!(aR)$);

  \node[label, above left=1pt] at (aUL) {$1'$};
  \node[label, below left=1pt] at (aLL) {$1$};
  \node[label, above right=1pt] at (aUR) {$2'$};
  \node[label, below right=1pt] at (aLR) {$2$};

  % Second diagram: the two outgoing fermions are interchanged.
  \coordinate (bL) at (1.45,0);
  \coordinate (bR) at (4.15,0);
  \coordinate (bUL) at (0.45,1.55);
  \coordinate (bLL) at (0.55,-1.45);
  \coordinate (bUR) at (5.15,1.55);
  \coordinate (bLR) at (5.05,-1.45);

  \draw[boson] (bL)--(bR);
  \draw[fermion] (bL)--(bUL);
  \draw[fermion] (bLL)--(bL);
  \draw[fermion] (bR)--(bUR);
  \draw[fermion] (bLR)--(bR);
  \draw[flow] ($(bL)!0.37!(bUL)$)--($(bL)!0.62!(bUL)$);
  \draw[flow] ($(bLL)!0.37!(bL)$)--($(bLL)!0.62!(bL)$);
  \draw[flow] ($(bR)!0.37!(bUR)$)--($(bR)!0.62!(bUR)$);
  \draw[flow] ($(bLR)!0.37!(bR)$)--($(bLR)!0.62!(bR)$);

  \node[label, above left=1pt] at (bUL) {$2'$};
  \node[label, below left=1pt] at (bLL) {$1$};
  \node[label, above right=1pt] at (bUR) {$1'$};
  \node[label, below right=1pt] at (bLR) {$2$};
\end{tikzpicture}
\end{center}

\caption{The connected second-order diagrams for fermion--fermion scattering
in a theory with interaction \eqref{eq:6.1.18}. Here straight lines represent
fermions; dotted lines are neutral bosons. There is a minus sign difference in
the contributions of these two diagrams, arising from an extra interchange of
fermion operators in the pairings represented by the second diagram.}
\label{fig:6.4}
\end{figure}
\endgroup

However, it must not be thought that all sign factors can be related in such a
simple way to the antisymmetry of the final or initial states, even in the
lowest order of perturbation theory. To illustrate this point, let's now
consider fermion--antifermion scattering, $1,2^c\to1'2'{}^c$, to second order
in the same interaction \eqref{eq:6.1.18}. The fermionic operators in the second-order
term in Eq. \eqref{eq:6.1.1} appear in the order:
\begin{equation}
a_{2'{}^c}a_{1'}\psi^\dagger(x)\psi(x)\psi^\dagger(y)\psi(y)
 a^\dagger_{1}a^\dagger_{2^c}.
\tag{6.1.22}
\label{eq:6.1.22}
\end{equation}
Here again there are two Feynman diagrams to this order, corresponding to the
pairings
\begin{equation}
[a_{2'{}^c}\psi(x)]\,[a_{1'}\psi^\dagger(x)]
 [\psi(y)a^\dagger_{1}]\,[\psi^\dagger(y)a^\dagger_{2^c}]
\tag{6.1.23}
\label{eq:6.1.23}
\end{equation}
and
\begin{equation}
[a_{2'{}^c}\psi(x)]\,[a_{1'}\psi^\dagger(y)]
 [\psi(y)a^\dagger_{1}]\,[\psi^\dagger(x)a^\dagger_{2^c}].
\tag{6.1.24}
\label{eq:6.1.24}
\end{equation}
(See Figure~\ref{fig:6.5}.) To go from \eqref{eq:6.1.22} to \eqref{eq:6.1.23} requires an even permutation
of fermionic operators (for instance, move $\psi(x)$ past two operators to the
left and move $\psi^\dagger(y)$ past two operators to the right) so there is no
extra minus sign in the contribution of the pairing \eqref{eq:6.1.23}. On the other
hand, to go from \eqref{eq:6.1.22} to \eqref{eq:6.1.24} requires an odd permutation of fermionic
operators (the same as for \eqref{eq:6.1.23}, \emph{plus} the interchange of
$\psi^\dagger(x)$ and $\psi^\dagger(y)$) so the contribution of this pairing
does come with an extra minus sign.\footnote{Actually, this sign is not wholly
unrelated to the requirements of Fermi statistics. The same field can destroy
a particle and create an antiparticle, so there is a relation, known as
`crossing symmetry', between processes in which initial particles or
antiparticles are exchanged with final antiparticles or particles. In
particular the amplitudes for the process $1,2^c\to1'2'{}^c$ are related to
those for the `crossed' process $1,2'\to1'2$; the two pairings \eqref{eq:6.1.23} and
\eqref{eq:6.1.24} just correspond to the two diagrams for this process, which differ by
an interchange of 1 and $2'$ (or $1'$ and 2), so the antisymmetry of the
scattering amplitude under interchange of initial (or final) particles
naturally requires a minus sign in the relative contribution of these two
pairings. However, crossing symmetry is not an ordinary symmetry (it involves
an analytic continuation in kinematic variables) and it is difficult to use it
with any precision for general processes.}

\begingroup
\renewcommand{\thefigure}{6.\arabic{figure}}
\setcounter{figure}{4}
\begin{figure}[htbp]
% Figure 6.5 diagram fragment. Caption remains in sec61b.tex.
% Source: physical PDF p. 293 (printed p. 270).
\begin{center}
\begin{tikzpicture}[
  fermion/.style={draw=black, line width=0.8pt},
  boson/.style={draw=black, line width=0.75pt,
    dash pattern=on 3pt off 2.5pt},
  flow/.style={-{Stealth[length=2.4mm,width=1.8mm]}, line width=0.8pt},
  label/.style={font=\small, inner sep=1pt}
]
  % First diagram: exchange channel.
  \coordinate (aL) at (-4.25,0);
  \coordinate (aR) at (-1.55,0);
  \coordinate (aUL) at (-5.25,1.55);
  \coordinate (aLL) at (-5.15,-1.45);
  \coordinate (aUR) at (-0.55,1.55);
  \coordinate (aLR) at (-0.65,-1.45);

  \draw[boson] (aL)--(aR);
  \draw[fermion] (aL)--(aUL);
  \draw[fermion] (aLL)--(aL);
  \draw[fermion] (aR)--(aUR);
  \draw[fermion] (aR)--(aLR);
  \draw[flow] ($(aL)!0.37!(aUL)$)--($(aL)!0.62!(aUL)$);
  \draw[flow] ($(aLL)!0.37!(aL)$)--($(aLL)!0.62!(aL)$);
  \draw[flow] ($(aUR)!0.37!(aR)$)--($(aUR)!0.62!(aR)$);
  \draw[flow] ($(aR)!0.37!(aLR)$)--($(aR)!0.62!(aLR)$);

  \node[label, above left=1pt] at (aUL) {$1'$};
  \node[label, below left=1pt] at (aLL) {$1$};
  \node[label, above right=1pt] at (aUR) {$2^{c'}$};
  \node[label, below right=1pt] at (aLR) {$2^c$};

  % Second diagram: annihilation channel, with vertical neutral-boson exchange.
  \coordinate (bT) at (3.25,0.95);
  \coordinate (bB) at (3.25,-0.95);
  \coordinate (bUL) at (2.15,2.15);
  \coordinate (bUR) at (4.35,2.15);
  \coordinate (bLL) at (2.15,-2.15);
  \coordinate (bLR) at (4.35,-2.15);

  \draw[boson] (bT)--(bB);
  \draw[fermion] (bT)--(bUL);
  \draw[fermion] (bT)--(bUR);
  \draw[fermion] (bLL)--(bB);
  \draw[fermion] (bB)--(bLR);
  \draw[flow] ($(bT)!0.37!(bUL)$)--($(bT)!0.62!(bUL)$);
  \draw[flow] ($(bUR)!0.37!(bT)$)--($(bUR)!0.62!(bT)$);
  \draw[flow] ($(bLL)!0.37!(bB)$)--($(bLL)!0.62!(bB)$);
  \draw[flow] ($(bB)!0.37!(bLR)$)--($(bB)!0.62!(bLR)$);

  \node[label, above left=1pt] at (bUL) {$1'$};
  \node[label, above right=1pt] at (bUR) {$2^{c'}$};
  \node[label, below left=1pt] at (bLL) {$1$};
  \node[label, below right=1pt] at (bLR) {$2^c$};
\end{tikzpicture}
\end{center}

\caption{The connected second-order diagrams for fermion--antifermion
scattering in a theory with interaction \eqref{eq:6.1.18}. Here straight lines
represent fermions or antifermions, depending on the arrow direction; dotted
lines are neutral bosons. There is again a minus sign difference in the
contributions of these two diagrams, arising from an extra interchange of
fermion operators in the pairings represented by the second diagram.}
\label{fig:6.5}
\end{figure}
\endgroup

Additional signs are encountered when we consider contributions of higher
order. In theories of the type considered here, in which the interactions of
fermions all take the form \eqref{eq:6.1.18}, the fermion lines in general Feynman
diagrams form either chains of lines that pass through the diagram with
arbitrary numbers of interactions with the boson fields, as in Figure~\ref{fig:6.6}, or
else fermionic \emph{loops}, like that shown in Figure~\ref{fig:6.7}. Consider the effect
of adding a fermionic loop with $M$ corners to the Feynman diagram for any
process. This corresponds to the pairing of fermionic operators
\begin{equation}
[\psi(x_1)\psi^\dagger(x_2)]
[\psi(x_2)\psi^\dagger(x_3)]\cdots
[\psi(x_M)\psi^\dagger(x_1)].
\tag{6.1.25}
\label{eq:6.1.25}
\end{equation}
On the other hand, these operators appear in Eq. \eqref{eq:6.1.1} in the order
\begin{equation}
\psi^\dagger(x_1)\psi(x_1)\psi^\dagger(x_2)\psi(x_2)\cdots
 \psi^\dagger(x_M)\psi(x_M).
\tag{6.1.26}
\label{eq:6.1.26}
\end{equation}
To go from \eqref{eq:6.1.26} to \eqref{eq:6.1.25} requires an odd permutation of fermionic
operators (move $\psi^\dagger(x_1)$ to the right past $2M-1$ operators) so the
contribution of each such fermionic loop is accompanied with a minus sign.

These rules yield the full $S$-matrix, including contributions from processes
in which various clusters of particles interact in widely separated regions
of spacetime. As discussed in Chapter 4, to calculate the part of the
$S$-matrix that excludes such contributions, we should include only
\emph{connected} Feynman diagrams. In particular, this excludes lines passing
clean through the diagram without interacting, which would yield the factors
\eqref{eq:6.1.13}.

To make the Feynman rules perfectly clear, we will calculate the low-order
contributions to the $S$-matrix for particle scattering in two different
theories.

\begingroup
\renewcommand{\thefigure}{6.\arabic{figure}}
\setcounter{figure}{5}
\begin{figure}[htbp]
% Figure 6.6 diagram fragment. Caption remains in sec61b.tex.
% Source: physical PDF p. 294 (printed p. 271).
\begin{center}
\begin{tikzpicture}[
  fermion/.style={draw=black, line width=0.8pt},
  boson/.style={draw=black, line width=0.75pt,
    dash pattern=on 3pt off 2.5pt},
  flow/.style={-{Stealth[length=2.4mm,width=1.8mm]}, line width=0.8pt},
  label/.style={font=\small, inner sep=1pt}
]
  % First diagram: the internal fermion line is vertical.
  \coordinate (aT) at (-3.25,0.95);
  \coordinate (aB) at (-3.25,-0.95);
  \coordinate (aUL) at (-4.45,2.15);
  \coordinate (aUR) at (-2.05,2.15);
  \coordinate (aLL) at (-4.35,-2.15);
  \coordinate (aLR) at (-2.05,-2.15);

  \draw[fermion] (aB)--(aT);
  \draw[fermion] (aT)--(aUL);
  \draw[fermion] (aLL)--(aB);
  \draw[boson] (aT)--(aUR);
  \draw[boson] (aB)--(aLR);
  \draw[flow] ($(aB)!0.38!(aT)$)--($(aB)!0.62!(aT)$);
  \draw[flow] ($(aT)!0.37!(aUL)$)--($(aT)!0.62!(aUL)$);
  \draw[flow] ($(aLL)!0.37!(aB)$)--($(aLL)!0.62!(aB)$);

  \node[label, above left=1pt] at (aUL) {$1'$};
  \node[label, above right=1pt] at (aUR) {$2'$};
  \node[label, below left=1pt] at (aLL) {$1$};
  \node[label, below right=1pt] at (aLR) {$2$};

  % Second diagram: the internal fermion line is horizontal.
  \coordinate (bL) at (1.15,0);
  \coordinate (bR) at (3.85,0);
  \coordinate (bUL) at (0.15,1.55);
  \coordinate (bLL) at (0.25,-1.45);
  \coordinate (bUR) at (4.85,1.55);
  \coordinate (bLR) at (4.75,-1.45);

  \draw[fermion] (bL)--(bR);
  \draw[fermion] (bLL)--(bL);
  \draw[fermion] (bR)--(bUR);
  \draw[boson] (bL)--(bUL);
  \draw[boson] (bR)--(bLR);
  \draw[flow] ($(bL)!0.38!(bR)$)--($(bL)!0.62!(bR)$);
  \draw[flow] ($(bLL)!0.37!(bL)$)--($(bLL)!0.62!(bL)$);
  \draw[flow] ($(bR)!0.37!(bUR)$)--($(bR)!0.62!(bUR)$);

  \node[label, above left=1pt] at (bUL) {$2'$};
  \node[label, below left=1pt] at (bLL) {$1$};
  \node[label, above right=1pt] at (bUR) {$1'$};
  \node[label, below right=1pt] at (bLR) {$2$};
\end{tikzpicture}
\end{center}

\caption{The connected second-order diagrams for boson--fermion scattering in
a theory with interaction \eqref{eq:6.1.18}. Straight lines are fermions;
dashed lines are neutral bosons.}
\label{fig:6.6}
\end{figure}
\endgroup

\begingroup
\renewcommand{\thefigure}{6.\arabic{figure}}
\setcounter{figure}{6}
\begin{figure}[htbp]
% Figure 6.7 diagram plate only; caption remains in sec61b.tex.
\begin{center}
\begin{tikzpicture}[
  line cap=round,
  line join=round,
  fermion/.style={draw, line width=0.75pt},
  boson/.style={draw, dashed, dash pattern=on 3pt off 3pt, line width=0.75pt}
]
  \coordinate (nw) at (-1.15,1.15);
  \coordinate (ne) at (1.15,1.15);
  \coordinate (se) at (1.15,-1.15);
  \coordinate (sw) at (-1.15,-1.15);

  \draw[fermion] (nw) -- (ne) -- (se) -- (sw) -- cycle;
  \draw[-{Latex[length=2.5mm,width=1.6mm]},line width=0.75pt]
    (-0.28,1.15) -- (0.30,1.15);
  \draw[-{Latex[length=2.5mm,width=1.6mm]},line width=0.75pt]
    (1.15,0.28) -- (1.15,-0.30);
  \draw[-{Latex[length=2.5mm,width=1.6mm]},line width=0.75pt]
    (0.28,-1.15) -- (-0.30,-1.15);
  \draw[-{Latex[length=2.5mm,width=1.6mm]},line width=0.75pt]
    (-1.15,-0.28) -- (-1.15,0.30);

  \draw[boson] (nw) -- (-2.35,2.05);
  \draw[boson] (ne) -- (2.35,2.05);
  \draw[boson] (se) -- (2.35,-2.05);
  \draw[boson] (sw) -- (-2.35,-2.05);
\end{tikzpicture}
\end{center}

\caption{The lowest-order connected diagram for boson--boson scattering in a
theory with interaction \eqref{eq:6.1.18}. Such fermion loop graphs yield an
extra minus sign, arising from permutations of the paired fermion fields.}
\label{fig:6.7}
\end{figure}
\endgroup

\begin{center}
\textit{Theory I}
\end{center}

Consider the theory of fermions and self-charge-conjugate bosons with
interaction \eqref{eq:6.1.18}. The lowest-order connected diagrams for fermion--boson
scattering are shown in Figure~\ref{fig:6.6}. Following the rules outlined in Figure~\ref{fig:6.1},
the corresponding $S$-matrix element is
\begin{align}
&S_{\mathbf p'_1\sigma'_1n'_1\,\mathbf p'_2\sigma'_2n'_2,
       \mathbf p_1\sigma_1n_1\,\mathbf p_2\sigma_2n_2}
\notag\\
&\quad=(2\pi)^{-6}\sum_{\ell'm'k'\ell m k}
 (-\ii)^2 g_{\ell'm'k'}g_{\ell m k}
 u_{\ell'}^*(\mathbf p'_1\sigma'_1n'_1)
 u_\ell(\mathbf p_1\sigma_1n_1)
\notag\\
&\qquad\cdot\int \dd^4x\int \dd^4y\,
 \bigl(-\ii\Delta_{m'm}(y-x)\bigr)
 e^{-\ii p'_1\cdot y}e^{\ii p_1\cdot x}
\notag\\
&\qquad\cdot\Bigl[
 e^{-\ii p'_2\cdot y}u_{k'}^*(\mathbf p'_2\sigma'_2n'_2)
 e^{\ii p_2\cdot x}u_k(\mathbf p_2\sigma_2n_2)
\notag\\
&\hspace{77mm}+
 e^{-\ii p'_2\cdot x}u_k^*(\mathbf p'_2\sigma'_2n'_2)
 e^{\ii p_2\cdot y}u_{k'}(\mathbf p_2\sigma_2n_2)
 \Bigr].
\tag{6.1.27}
\label{eq:6.1.27}
\end{align}
(The labels 1 and 2 are used here for fermions and bosons, respectively.) For
fermion--fermion scattering there are also two second-order diagrams, shown in
Figure~\ref{fig:6.4}. They yield the $S$-matrix element
\begin{align}
&S_{\mathbf p'_1\sigma'_1n'_1\,\mathbf p'_2\sigma'_2n'_2,
       \mathbf p_1\sigma_1n_1\,\mathbf p_2\sigma_2n_2}
\notag\\
&\quad=(2\pi)^{-6}\sum_{k'\ell'm'k\ell m}
 (-\ii)^2g_{m'\ell k'}g_{\ell'mk}
 u_{m'}^*(\mathbf p'_2\sigma'_2n'_2)
 u_{\ell'}^*(\mathbf p'_1\sigma'_1n'_1)
 u_m(\mathbf p_2\sigma_2n_2)u_\ell(\mathbf p_1\sigma_1n_1)
\notag\\
&\qquad\cdot\int \dd^4x\int \dd^4y\,
 e^{-\ii p'_2\cdot x}e^{-\ii p'_1\cdot y}
 e^{\ii p_2\cdot x}e^{\ii p_1\cdot y}
 \bigl(-\ii\Delta_{k'k}(x-y)\bigr)
 -[1'\rightleftarrows2'].
\tag{6.1.28}
\label{eq:6.1.28}
\end{align}
with the last term indicating subtraction of the preceding term with
interchange of particles $1'$ and $2'$ (or equivalently 1 and 2). There are no
second-order graphs for boson--boson scattering in this theory; the
lowest-order graphs are of fourth order, such as that shown in Figure~\ref{fig:6.7}. More
specific examples of formulas like Eqs. \eqref{eq:6.1.27} and \eqref{eq:6.1.28} will be given in
Section 6.3, after we have had a chance to evaluate the propagators and go over
to momentum space.

\emph{In the interaction \eqref{eq:6.1.18} of the preceding example, the three fields
are all different.} It is instructive also to look at an example with a
trilinear interaction in which all three fields are the same, or at least enter
into the interaction in a symmetric way.

\begin{center}
\textit{Theory II}
\end{center}

Now take the interaction density to be a sum of terms that are trilinear in a
set of \emph{real} bosonic fields $\phi_\ell(x)$:
\begin{equation}
\mathcal H_I(x)=\frac{1}{3!}\sum_{\ell mn}g_{\ell mn}
 \phi_\ell(x)\phi_m(x)\phi_n(x)
\tag{6.1.29}
\label{eq:6.1.29}
\end{equation}
with $g_{\ell mn}$ a real totally symmetric coupling coefficient. Suppose we
want to consider a scattering process $1,2\to1'2'$ to second order in this
interaction. Each of the two vertices must have two of the four external lines
attached to it. (The only other possibility is that one of the external lines
is attached to one vertex and three to the other vertex, but the vertex with
three external lines attached to it would have no remaining lines to connect
it to the other vertex so this would be a disconnected contribution.) The
additional line required at each vertex must then just serve to connect the two
vertices to each other. There are three graphs of this type, differing in
whether the other external line that is attached to the same vertex as line 1
is line 2 or $1'$ or $2'$. (See Figure~\ref{fig:6.8}.) Following the rules given above,
the contribution to the $S$-matrix from those three diagrams is

\begingroup
\renewcommand{\thefigure}{6.\arabic{figure}}
\setcounter{figure}{7}
\begin{figure}[htbp]
% Figure 6.8 diagram plate only; caption remains in sec61b.tex.
\begin{center}
\begin{tikzpicture}[
  line cap=round,
  line join=round,
  scalar/.style={draw, dashed, dash pattern=on 3pt off 3pt, line width=0.75pt},
  every node/.style={font=\small}
]
  % First channel: vertical internal line.
  \coordinate (a1) at (-5.0,0.95);
  \coordinate (a2) at (-5.0,-0.95);
  \draw[scalar] (a1) -- (a2);
  \draw[scalar] (a1) -- (-5.75,2.05);
  \draw[scalar] (a1) -- (-4.25,2.05);
  \draw[scalar] (a2) -- (-5.75,-2.05);
  \draw[scalar] (a2) -- (-4.25,-2.05);
  \node[above] at (-5.75,2.05) {$1'$};
  \node[above] at (-4.25,2.05) {$2'$};
  \node[below] at (-5.75,-2.05) {$1$};
  \node[below] at (-4.25,-2.05) {$2$};

  % Second channel: 1 and 1' meet at the left vertex.
  \coordinate (b1) at (-1.55,0);
  \coordinate (b2) at (0.05,0);
  \draw[scalar] (b1) -- (b2);
  \draw[scalar] (b1) -- (-2.55,0.85);
  \draw[scalar] (b1) -- (-2.55,-0.85);
  \draw[scalar] (b2) -- (1.05,0.85);
  \draw[scalar] (b2) -- (1.05,-0.85);
  \node[above] at (-2.55,0.85) {$1'$};
  \node[below] at (-2.55,-0.85) {$1$};
  \node[above] at (1.05,0.85) {$2'$};
  \node[below] at (1.05,-0.85) {$2$};

  % Third channel: 1 and 2' meet at the left vertex.
  \coordinate (c1) at (3.15,0);
  \coordinate (c2) at (4.75,0);
  \draw[scalar] (c1) -- (c2);
  \draw[scalar] (c1) -- (2.15,0.85);
  \draw[scalar] (c1) -- (2.15,-0.85);
  \draw[scalar] (c2) -- (5.75,0.85);
  \draw[scalar] (c2) -- (5.75,-0.85);
  \node[above] at (2.15,0.85) {$2'$};
  \node[below] at (2.15,-0.85) {$1$};
  \node[above] at (5.75,0.85) {$1'$};
  \node[below] at (5.75,-0.85) {$2$};
\end{tikzpicture}
\end{center}

\caption{The connected-second order diagrams for boson--boson scattering in a
theory with interaction \eqref{eq:6.1.29}.}
\label{fig:6.8}
\end{figure}
\endgroup

\begin{align}
&S_{\mathbf p'_1\sigma'_1n'_1\,\mathbf p'_2\sigma'_2n'_2,
       \mathbf p_1\sigma_1n_1\,\mathbf p_2\sigma_2n_2}
\notag\\
&\quad=(-\ii)^2(2\pi)^{-6}
 \sum_{\ell\ell'\ell''mm'm''}
 g_{\ell\ell'\ell''}g_{mm'm''}
 \int \dd^4x\int \dd^4y\,
 \bigl(-\ii\Delta_{\ell''m''}(x,y)\bigr)
\notag\\
&\qquad\cdot\Bigl[
 u_\ell^*(\mathbf p'_1\sigma'_1n'_1)e^{-\ii p'_1\cdot x}
 u_{\ell'}^*(\mathbf p'_2\sigma'_2n'_2)e^{-\ii p'_2\cdot x}
 u_m(\mathbf p_1\sigma_1n_1)e^{\ii p_1\cdot y}
 u_{m'}(\mathbf p_2\sigma_2n_2)e^{\ii p_2\cdot y}
\notag\\
&\qquad\quad+
 u_\ell^*(\mathbf p'_1\sigma'_1n'_1)e^{-\ii p'_1\cdot x}
 u_{\ell'}(\mathbf p_1\sigma_1n_1)e^{\ii p_1\cdot x}
 u_{m'}^*(\mathbf p'_2\sigma'_2n'_2)e^{-\ii p'_2\cdot y}
 u_m(\mathbf p_2\sigma_2n_2)e^{\ii p_2\cdot y}
\notag\\
&\qquad\quad+
 u_\ell^*(\mathbf p'_2\sigma'_2n'_2)e^{-\ii p'_2\cdot x}
 u_{\ell'}(\mathbf p_1\sigma_1n_1)e^{\ii p_1\cdot x}
 u_{m'}^*(\mathbf p'_1\sigma'_1n'_1)e^{-\ii p'_1\cdot y}
 u_m(\mathbf p_2\sigma_2n_2)e^{\ii p_2\cdot y}
 \Bigr].
\tag{6.1.30}
\label{eq:6.1.30}
\end{align}
To be even more specific, if the bosons in this theory are spinless particles
of a single species, then we write the interaction \eqref{eq:6.1.29} in the form
\begin{equation}
\mathcal H_I=g\phi^3/3!
\tag{6.1.31}
\label{eq:6.1.31}
\end{equation}
and the $S$-matrix element \eqref{eq:6.1.30} for scalar--scalar scattering is
\begin{align*}
S_{\mathbf p'_1\,\mathbf p'_2,\,\mathbf p_1\,\mathbf p_2}
={}&\frac{ig^2}{(2\pi)^6\sqrt{16E'_1E'_2E_1E_2}}
 \int \dd^4x\int \dd^4y\,\Delta_F(x-y)
\notag\\
&\cdot\Bigl[
 \exp\bigl(-\ii(p'_1+p'_2)\cdot x\bigr)
 \exp\bigl(\ii(p_1+p_2)\cdot y\bigr)
\notag\\
&\qquad+
 \exp\bigl(\ii(p_1-p'_1)\cdot x\bigr)
 \exp\bigl(\ii(p_2-p'_2)\cdot y\bigr)
\notag\\
&\qquad+
 \exp\bigl(\ii(p_1-p'_2)\cdot x\bigr)
 \exp\bigl(\ii(p_2-p'_1)\cdot y\bigr)
 \Bigr],
\end{align*}
where $\Delta_F(x-y)$ is the scalar field propagator, calculated in the next
section. There are no terms of third order in $\mathcal H_I(x)$, or of any odd
order in $\mathcal H_I(x)$.
